\documentclass[aspectratio=169,11pt]{beamer}
% Shared CMPSC 480 Beamer preamble.
\usepackage[T1]{fontenc}
\usepackage[utf8]{inputenc}
\usepackage{lmodern}
\usepackage{amsmath,amssymb,mathtools}
\usepackage{booktabs,array,tabularx}
\usepackage{xcolor}
\usepackage{listings}
\usepackage{tikz}
\usetikzlibrary{arrows.meta,positioning,shapes.geometric}

% Ignore only negligible TeX box diagnostics that are below visual tolerance.
% Larger layout defects still appear in the build log and in rendered QA.
\vfuzz=3pt
\hbadness=2000

\definecolor{courseink}{HTML}{111111}
\definecolor{coursegray}{HTML}{EDEDED}
\definecolor{courserule}{HTML}{B8BCC4}
\definecolor{courseblue}{HTML}{3D8DFF}
\definecolor{courselightblue}{HTML}{D0EDFA}
\definecolor{coursemuted}{HTML}{5C6470}
\definecolor{coursegreen}{HTML}{0B7A53}
\definecolor{coursered}{HTML}{B3261E}

\usetheme{default}
\usefonttheme{professionalfonts}
\setbeamertemplate{navigation symbols}{}
\setbeamersize{text margin left=0.65cm,text margin right=0.65cm}

\setbeamercolor{normal text}{fg=courseink,bg=white}
\setbeamercolor{frametitle}{fg=courseink,bg=white}
\setbeamercolor{title}{fg=courseink,bg=white}
\setbeamercolor{subtitle}{fg=coursemuted,bg=white}
\setbeamercolor{structure}{fg=courseblue}
\setbeamercolor{alerted text}{fg=coursered}
\setbeamercolor{block title}{fg=courseink,bg=courselightblue}
\setbeamercolor{block body}{fg=courseink,bg=coursegray}
\setbeamercolor{example text}{fg=coursegreen}

\setbeamerfont{title}{size=\fontsize{25}{29}\selectfont,series=\bfseries}
\setbeamerfont{subtitle}{size=\fontsize{13}{16}\selectfont}
\setbeamerfont{frametitle}{size=\fontsize{19}{22}\selectfont,series=\bfseries}
\setbeamerfont{normal text}{size=\fontsize{11}{14}\selectfont}
\setbeamerfont{footline}{size=\fontsize{7.5}{9}\selectfont}

\setbeamertemplate{frametitle}{%
  \nointerlineskip
  % Use content-driven height so a long assessment title can wrap without
  % being clipped by a fixed-height title bar.
  \begin{beamercolorbox}[wd=\paperwidth,sep=0.17cm,leftskip=0.48cm,rightskip=0.48cm]{frametitle}
    \insertframetitle
  \end{beamercolorbox}
  \vspace{-0.08cm}
  {\color{courserule}\rule{\textwidth}{0.35pt}}
}

\setbeamertemplate{footline}{%
  \leavevmode
  \begin{beamercolorbox}[wd=\paperwidth,ht=2.6ex,dp=1.1ex,leftskip=.65cm,rightskip=.65cm]{author in head/foot}
    \color{coursemuted}\insertshorttitle\hfill\insertframenumber/\inserttotalframenumber
  \end{beamercolorbox}
}

\setbeamertemplate{itemize item}{\color{courseblue}\small$\blacktriangleright$}
\setbeamertemplate{itemize subitem}{\color{courseblue}\scriptsize$\bullet$}
\setbeamertemplate{enumerate item}{\color{courseblue}\insertenumlabel.}

\lstdefinestyle{coursepython}{
  language=Python,
  basicstyle=\ttfamily\fontsize{8.5}{10}\selectfont,
  keywordstyle=\color{courseblue}\bfseries,
  commentstyle=\color{coursemuted}\itshape,
  stringstyle=\color{coursegreen},
  showstringspaces=false,
  columns=fullflexible,
  keepspaces=true,
  frame=single,
  rulecolor=\color{courserule},
  backgroundcolor=\color{coursegray},
  breaklines=true,
  tabsize=4,
  xleftmargin=0.15cm,
  xrightmargin=0.15cm,
  framexleftmargin=0.15cm,
  framexrightmargin=0.15cm
}
\lstset{style=coursepython}

\newcommand{\points}[1]{\hfill{\color{courseblue}\textbf{[#1 points]}}}
\newcommand{\deliverable}[1]{\textbf{\color{courseblue}#1}}
\newcommand{\code}[1]{\texttt{#1}}
\newcommand{\keyidea}[1]{%
  \begin{beamercolorbox}[rounded=false,sep=0.55em,wd=\linewidth]{block body}
    \textbf{Key idea.} #1
  \end{beamercolorbox}
}
\newcommand{\answerlabel}{\textbf{\color{coursegreen}Answer.}\ }
\newcommand{\gradinglabel}{\textbf{\color{courseblue}Grading guidance.}\ }

\newenvironment{questionblock}[1]
  {\begin{block}{#1}}
  {\end{block}}

\newenvironment{solutionblock}[1]
  {\begin{block}{#1}}
  {\end{block}}

\AtBeginSection[]{
  \begin{frame}
    \vfill
    {\usebeamerfont{title}\color{courseink}\insertsectionhead\par}
    \vspace{0.4cm}
    {\color{courseblue}\rule{0.32\paperwidth}{3pt}}
    \vfill
  \end{frame}
}


% CMPSC480_TOTAL_POINTS: 10
% CMPSC480_ITEM_IDS: 1,2,3,4
\title[Week 6 Assignment Questions]{Week 6 Assignment: Linear Q-Function Approximation}
\subtitle{Feature design, semi-gradient updates, and stability evidence}
\author{CMPSC 480 - Student Question Deck}
\date{}

\begin{document}


\begin{frame}
  \titlepage
\end{frame}

\begin{frame}{Assessment overview}
\begin{columns}[T,onlytextwidth]
\begin{column}{0.62\textwidth}
\Large Replace a Q-table with a fixed-size linear approximator and verify every update.

\vspace{0.35cm}
\normalsize
Coverage: Week 6: function approximation, semi-gradients, and the deadly triad
\end{column}
\begin{column}{0.33\textwidth}
\begin{block}{Points}10 total\end{block}
\begin{block}{Expected time}6-8 hours\end{block}
\begin{block}{Submission}Repository plus experiment memo\end{block}
\end{column}
\end{columns}
\end{frame}

\begin{frame}{Learning objectives}
\begin{itemize}
  \item Implement a fixed-length state-action feature map.
  \item Compute Q values as vectorized dot products.
  \item Apply terminal-aware semi-gradient Q-learning.
  \item Demonstrate the one-hot tabular correspondence.
  \item Monitor numerical stability across controlled runs.
\end{itemize}
\end{frame}

\begin{frame}{Requirements checklist}
\small
\begin{itemize}
  \item Use Python 3.11 and NumPy; do not use a library that implements the assessed algorithm.
  \item Keep data, model, optimization, and evaluation responsibilities behind explicit interfaces.
  \item Validate shapes, ranges, finite values, empty inputs, and terminal or degenerate cases.
  \item Use deterministic seeds and record every configuration used for reported evidence.
  \item Include focused tests and a README with exact reproduction commands.
\end{itemize}
\end{frame}

\begin{frame}{Task 1 - Features and dimensions (2 points)}
\small
Design a feature map whose length is independent of the number of raw states.

\vspace{0.2cm}
\begin{enumerate}
\item[\textbf{a.}] Define at least four state-action features and justify their scales. \hfill{\color{courseblue}\textbf{[1 points]}}
\item[\textbf{b.}] Prove by an executable assertion that parameter size does not grow when the environment map grows. \hfill{\color{courseblue}\textbf{[1 points]}}
\end{enumerate}
\end{frame}

\begin{frame}{Task 2 - Linear values and policy (2 points)}
\small
Implement Q-hat and epsilon-greedy action selection over legal actions.

\vspace{0.2cm}
\begin{enumerate}
\item[\textbf{a.}] Compute Q-hat(s,a;w) and test a known vector by hand. \hfill{\color{courseblue}\textbf{[1 points]}}
\item[\textbf{b.}] Specify training and evaluation tie behavior. \hfill{\color{courseblue}\textbf{[1 points]}}
\end{enumerate}
\end{frame}

\begin{frame}{Task 3 - Semi-gradient update (4 points)}
\small
Implement the squared-TD-error semi-gradient rule with correct terminal masking.

\vspace{0.2cm}
\begin{enumerate}
\item[\textbf{a.}] Derive the update from one-half times squared TD error. \hfill{\color{courseblue}\textbf{[2 points]}}
\item[\textbf{b.}] Compute an update for w=[0,0], f=[1,2], reward=1, gamma=0.5, max-next-Q=4, and alpha=0.1. \hfill{\color{courseblue}\textbf{[1 points]}}
\item[\textbf{c.}] State and test the terminal target. \hfill{\color{courseblue}\textbf{[1 points]}}
\end{enumerate}
\end{frame}

\begin{frame}{Task 4 - Experiment and stability (2 points)}
\small
Compare feature sets and report evidence beyond episodic return.

\vspace{0.2cm}
\begin{enumerate}
\item[\textbf{a.}] Run one feature ablation and compare with a tabular small-state oracle. \hfill{\color{courseblue}\textbf{[1 points]}}
\item[\textbf{b.}] Define an instability gate using at least three logged quantities. \hfill{\color{courseblue}\textbf{[1 points]}}
\end{enumerate}
\end{frame}

\begin{frame}{Edge cases and defensive programming}
\small
\begin{itemize}
  \item Zero feature vector.
  \item Terminal next state.
  \item Feature dimension mismatch.
  \item All legal actions tie.
  \item Weight or TD error becomes nonfinite.
\end{itemize}
\end{frame}

\begin{frame}{Submission instructions}
\small
\begin{itemize}
  \item Submit source, tests, README, configuration, and the requested report.
  \item Provide one command that reproduces the primary result from a clean environment.
  \item Exclude credentials, generated caches, and unlicensed data.
\end{itemize}
\vspace{0.2cm}
\alert{Do not submit credentials, generated caches, or unapproved libraries.}
\end{frame}

\begin{frame}{Grading rubric}
\scriptsize
\begin{tabularx}{\textwidth}{>{\bfseries}p{0.28\textwidth} p{0.08\textwidth} X}
\toprule
Category & Points & Required evidence \\
\midrule
Features and dimensions & 2 & Feature implementation, dimensionality assertion, and feature table. \\
Linear values and policy & 2 & Vectorized q\_value and deterministic evaluation tests. \\
Semi-gradient update & 4 & Update method, derivation, and hand-computed tests. \\
Experiment and stability & 2 & Ablation table, stability plots or summaries, and defensive tests. \\
\bottomrule
\end{tabularx}
\end{frame}

\begin{frame}{Reflection prompt}
In 150-250 words:

Identify one feature that generalized useful information and one feature that created coupling risk. Cite an ablation or trace.
\end{frame}

\end{document}
